% reporte.tex — Benchmarking de MongoDB frente a PostgreSQL JSONB bajo contención
% Bases de Datos en la Nube · IT2iD · Universidad Politécnica de Chiapas · 2026
%
% CÓMO SE COMPILA
%   Overleaf: subir la carpeta report/ completa y compilar con pdfLaTeX.
%   Local:    latexmk -pdf reporte.tex   (necesita pgfplots, de TeX Live)
%
% NINGÚN NÚMERO SE ESCRIBE AQUÍ A MANO (criterio de PLAN F9). Todo sale de:
%   cifras/cifras.tex   macros con las cifras del CSV   (\cifraFilasCsv, …)
%   tablas/*.tex        cuerpos de tabla ya formateados
%   datos/*.dat         series que pgfplots dibuja
% Los tres los genera  scripts/build_report_data.py  desde results/measurements.csv
% y results/raw/. Si cambian los datos, se regenera y el reporte queda al día.

\documentclass[11pt,a4paper]{article}

\usepackage[spanish,es-tabla]{babel}
\usepackage[utf8]{inputenc}
\usepackage[T1]{fontenc}
\usepackage{lmodern}
\usepackage[margin=2.5cm]{geometry}
\usepackage{booktabs}
\usepackage{longtable}
\usepackage{tabularx}
\usepackage{graphicx}
\usepackage{hyperref}
\usepackage{xcolor}
\usepackage{siunitx}
\usepackage{ifthen}   % \ifthenelse en las macros de las figuras
\sisetup{group-separator={\,}, group-minimum-digits=4, detect-weight}
\usepackage{pgfplots}
\usepackage{pgfplotstable}
\pgfplotsset{compat=1.18}

% Las figuras leen los .dat de report/datos/
\pgfplotsset{
  table/search path={datos},
  % Un `nan` en un .dat es un punto que NO existe: en Mongo · C · Q2 · c=50 TODAS
  % las consultas de la ventana son primeras consultas, así que no hay P99 "en
  % régimen". Con unbounded coords=jump pgfplots corta la línea ahí, en vez de
  % fallar o de unir dos puntos como si el punto intermedio existiera.
  every axis/.append style={
    grid=major, grid style={gray!25}, tick label style={font=\small},
    label style={font=\small}, legend style={font=\small, draw=gray!40},
    width=0.92\linewidth, height=6cm, unbounded coords=jump,
  },
}
\pgfplotstableset{search path={datos}}

\hypersetup{colorlinks=true, linkcolor=blue!50!black, urlcolor=blue!50!black,
            pdftitle={Benchmarking de MongoDB frente a PostgreSQL JSONB bajo contención}}

% Notas al pie de las tablas. R-29.2 (decisión de Alix): la celda colapsada lleva
% nota al pie en TODA tabla donde aparezca, no una sola vez en el texto.
\newcommand{\tnotaColapsada}{\textsuperscript{\dag}}
\newcommand{\tnotaCensurada}{\textsuperscript{\ddag}}
\newcommand{\tnotaBajaDensidad}{\textsuperscript{*}}
\newcommand{\leyendaNotas}{%
  \footnotesize
  \tnotaColapsada{}~El motor estaba \textbf{colapsado} durante la ventana: sus
  errores son conexiones rechazadas, no consultas ejecutadas, y su
  \texttt{max\_rss\_mb} \textbf{no es eficiencia} (R-29.2).\par
  \tnotaCensurada{}~Percentiles invalidados por censura por la derecha: la
  fracción de \emph{timeouts} supera el límite de A4.\par
  \tnotaBajaDensidad{}~Baja densidad muestral: menos de 20 consultas completadas
  en la ventana (A4).\par
}

% generado por scripts/build_report_data.py — no editar a mano
\newcommand{\cifraFilasCsv}{108}
\newcommand{\cifraCeldas}{36}
\newcommand{\cifraFilasOk}{78}
\newcommand{\cifraFilasSpill}{18}
\newcommand{\cifraFilasError}{7}
\newcommand{\cifraFilasOom}{5}
\newcommand{\cifraOomEnVentana}{2}
\newcommand{\cifraOomEnCalentamiento}{3}
\newcommand{\cifraFilasConSpill}{24}
\newcommand{\cifraTempTotalTB}{0.58}
\newcommand{\cifraCeldasMedibles}{32}
\newcommand{\cifraCeldasCensuradas}{4}
\newcommand{\cifraDescartadas}{0}


\title{\textbf{Benchmarking de MongoDB frente a PostgreSQL JSONB\\bajo contención de CPU y memoria}}
\author{Alix Montesinos\\\small Bases de Datos en la Nube · IT2iD\\\small Universidad Politécnica de Chiapas}
\date{Noviembre de 2026}

\begin{document}
\maketitle

\begin{abstract}
\noindent
Se comparan MongoDB y PostgreSQL con \texttt{JSONB} sobre el mismo conjunto de
\num{500000} documentos sintéticos, con dos consultas de agregación y tres
perfiles de recursos aplicados con \emph{cgroups} v2, a tres niveles de
concurrencia y con tres repeticiones por celda: \cifraFilasCsv{} mediciones en
\cifraCeldas{} celdas. El experimento mide lo que ocurre de verdad bajo
contención, incluidos los casos en los que un motor deja de responder, y declara
las limitaciones del entorno virtualizado en el que se ejecutó.
\end{abstract}

\tableofcontents
\newpage

% 01-entorno.tex — Sección 1 del manual (playbook.pdf §8.3): especificaciones del
% host y del entorno Docker. Incluye las limitaciones del entorno virtualizado,
% que van declaradas aquí y no escondidas (SPEC §6, L1–L7).

\section{Especificaciones del host y del entorno Docker}
\label{sec:entorno}

\subsection{Host y máquina virtual}

Las mediciones se ejecutaron en un portátil con Windows~11, WSL2 y Docker
Desktop. El repositorio y los volúmenes viven en el sistema de archivos de WSL2,
nunca bajo \texttt{/mnt/c/}, para no medir la capa de traducción de Windows.

\begin{table}[htbp]
  \centering
  \caption{Host y máquina virtual. Capturado por \texttt{scripts/capture\_env.py}
           en cada repetición; evidencia en \texttt{results/env/}.}
  \label{tab:host}
  \small
  \begin{tabularx}{\linewidth}{lX}
    \toprule
    Procesador & AMD Ryzen~5 PRO 5675U: 6 núcleos físicos, 12 hilos \\
    Memoria del host & 14.8~GiB \\
    Sistema & Windows~11 + WSL2 (distribución \texttt{Ubuntu}) + Docker Desktop \\
    VM de WSL2 & 12 CPUs lógicas y 7.159~GiB visibles \\
    \emph{cgroups} & v2, confirmado leyendo el cgroup desde dentro del contenedor \\
    Reparto de CPUs & motor en las lógicas 0--7; runner y telemetría en 8--11, sin
                      compartir núcleo físico (decisión F5-c, topología medida) \\
    \bottomrule
  \end{tabularx}
\end{table}

\subsection{Motores, fijados por digest}

Las imágenes no se referencian por etiqueta sino por \emph{digest}: una etiqueta
como \texttt{mongo:8.2} es móvil y un \texttt{docker pull} posterior podría traer
otra versión, lo que rompería la reproducibilidad (R-33).

\begin{table}[htbp]
  \centering
  \caption{Motores y telemetría. Digests completos en \texttt{docker/images.env} y
           en \texttt{results/env/f1/}.}
  \label{tab:imagenes}
  \small
  \begin{tabularx}{\linewidth}{llX}
    \toprule
    Componente & Versión & Digest (abreviado) \\
    \midrule
    MongoDB & 8.2.12 (wiredTiger) & \texttt{sha256:e0ce8c35\ldots{}4bd3} \\
    PostgreSQL & 18.6 & \texttt{sha256:3a426452\ldots{}4b81} \\
    Prometheus & v3.15.1 & \texttt{sha256:fa330227\ldots{}91a2} \\
    cAdvisor & v0.60.6 & \texttt{sha256:b8e7d109\ldots{}d07a} \\
    node\_exporter & v1.12.1 & \texttt{sha256:1b4e4438\ldots{}1be0} \\
    \bottomrule
  \end{tabularx}
\end{table}

\subsection{Perfiles de recursos}

Los tres perfiles del manual se aplican con \emph{cgroups} v2 y se
\textbf{verifican leyendo el cgroup} antes de cada repetición, no se dan por
puestos (\texttt{scripts/verify\_limits.sh}).

\begin{table}[htbp]
  \centering
  \caption{Perfiles y límites reales observados en el cgroup. Los seis
           (3 perfiles $\times$ 2 motores) se verificaron en F1.}
  \label{tab:perfiles}
  \small
  \begin{tabular}{llll}
    \toprule
    Perfil & CPU & Memoria & \texttt{memory.swap.max} \\
    \midrule
    \texttt{PROFILE\_A\_BASELINE} & 4.0 & 4096~m & 0 \\
    \texttt{PROFILE\_B\_MEDIUM}   & 2.0 & 2048~m & 0 \\
    \texttt{PROFILE\_C\_SEVERE}   & 1.0 & 1024~m & 0 \\
    \bottomrule
  \end{tabular}
\end{table}

\texttt{memory.swap.max = 0} en los seis casos: cuando un motor agota su límite
de memoria no tiene \emph{swap} donde escapar, y el kernel lo mata. Esa es la
garantía que sustituye a \texttt{mem\_swappiness}, que \textbf{no existe en
cgroups~v2}: el manual lo pide y se declara en el \emph{compose}, pero Docker lo
descarta con un aviso, así que no es un fallo de la configuración (SPEC~A7, V8).

\subsection{Limitaciones del entorno, declaradas}
\label{sec:limitaciones-entorno}

El manual supone un host dedicado. Este no lo es, y conviene decirlo antes de los
resultados y no después:

\begin{enumerate}
  \item \textbf{La limpieza de caché no llega hasta el disco.}
        \texttt{drop\_caches} vacía la caché de páginas de la VM de WSL2, pero
        \textbf{no la de Windows sobre el \texttt{.vhdx}} donde viven los
        volúmenes. Lo alcanzable es un estado de partida \emph{repetible}, no
        frío. Se midió el efecto en cada repetición (entre 0.7 y 1.8~GB
        liberados), no solo el código de salida del comando.
  \item \textbf{El perfil A no es un baseline holgado.} La VM tiene 7.159~GiB en
        total; con 4~GiB para el motor quedan unos 3~GiB para el runner, la
        telemetría y la caché de páginas. Se midió: en el peor caso quedaron
        3.2~GiB libres, sin \emph{swap}, pero con 120 de 300~s en espera por
        memoria (F7.4).
  \item \textbf{Las 4 CPUs del perfil A no son 4 núcleos físicos.} La VM expone 12
        lógicas sobre 6 físicos; el reparto 0--7 para el motor usa 4 núcleos
        físicos completos, medidos con \texttt{thread\_siblings\_list}.
  \item \textbf{La frecuencia del procesador baja bajo carga sostenida.} Medido
        durante la matriz: el rendimiento del procesador osciló entre el 86~\% y
        el \textbf{60.8~\%} de la frecuencia nominal (de 2296 a 1396~MHz). No se
        puede leer desde WSL2 —no hay \texttt{cpufreq}— y se midió desde Windows
        con el contador «\% de rendimiento del procesador». Es la causa probable
        de la mayor de las dispersiones observadas (sección~\ref{sec:resultados}).
  \item \textbf{cAdvisor observa la VM, no el equipo.} Las curvas de telemetría
        describen el contenedor dentro de WSL2.
\end{enumerate}

\subsection{Consecuencia para la limpieza de caché: permiso de \emph{root}}

El flush por repetición necesita \emph{root}. Para que una corrida de horas no se
detuviera a pedir contraseña se instaló una regla de \texttt{sudo} acotada a
\textbf{un único ejecutable} (\texttt{/usr/local/sbin/bench-drop-caches}),
propiedad de \texttt{root}, que no acepta parámetros y solo hace \texttt{sync} y
\texttt{drop\_caches}. Vive fuera del repositorio a propósito: si estuviera dentro,
donde el usuario puede editarlo, la regla equivaldría a \emph{root} sin
contraseña. El resto de \texttt{sudo} sigue pidiéndola, y la regla se deshace con
un script. Cada repetición registra por qué vía se limpió la caché
(\texttt{"via"} en el archivo crudo). Diseño y reversión en
\texttt{docs/V7\_SUDOERS.md}.

% 02-metodo.tex — Sección 2: método de medición y verificación del baseline.
% Responde las preguntas 1, 2 y 5 del revisor (diapo 11/13) y declara las
% decisiones de SPEC §5 en vez de esconderlas (criterio de PLAN F9).

\section{Método de medición y verificación del baseline}
\label{sec:metodo}

\subsection{El conjunto de datos es el mismo para los dos motores}
\label{sec:baseline}

Los \num{500000} documentos se generan con
\texttt{generator/synthetic\_data\_generator.py} con \textbf{semilla fija 42}, de
modo que el archivo es reproducible byte a byte. Su huella SHA-256 empieza por
\texttt{0cd7fadb} y termina en \texttt{0414}, y queda registrada en
\texttt{results/dataset\_stats.json} junto con las distribuciones obtenidas.

Que los dos motores contengan lo mismo no se supone: se verificó en tres
direcciones (archivo~$\to$~motor, motor~$\to$~archivo y motor~$\to$~motor) con
\texttt{scripts/verify\_gate.py}, comparando los documentos exportados de cada
motor contra la fuente en \emph{streaming}. Dos decisiones del esquema son las
que hacen que la comparación sea válida:

\begin{itemize}
  \item \texttt{payment.final\_total} se almacena como \textbf{número} JSON, no
        como cadena. Si fuera cadena, MongoDB ordenaría lexicográficamente
        (\texttt{"9.5" > "100.0"}) y PostgreSQL numéricamente: los Top-1000 no
        coincidirían y la comparación quedaría invalidada \emph{en silencio}.
  \item \texttt{timestamp} se almacena como cadena ISO-8601 en los dos motores.
        Como \texttt{Date} en MongoDB, el \texttt{\$gte} contra un literal de
        texto no haría \emph{match}.
\end{itemize}

\subsection{Las consultas se leen del disco, no están dentro del código}

Las cuatro consultas (\texttt{Q1} y \texttt{Q2}, para cada motor) se copiaron
\textbf{literalmente} del manual a \texttt{queries/} y el runner las lee de ahí en
cada ejecución, con una comprobación de integridad al arrancar. El revisor puede
compararlas con el PDF sin leer código.

Su equivalencia se verificó antes de medir: \texttt{scripts/check\_equivalence.py}
calcula el resultado de referencia en Python con aritmética exacta
(\texttt{Decimal}) y lo compara con lo que devuelve cada motor. Coinciden en todas
las variantes. El orden \emph{dentro} del Top-1000 de \texttt{Q2} no es
determinista —hay valores repetidos y la consulta no tiene desempate—, así que se
compara como conjunto.

\subsection{Qué es una repetición}

\begin{enumerate}
  \item Se baja el motor que estuviera arriba. \textbf{Solo un motor corre a la
        vez}: con los dos arriba competirían por la caché de páginas de la VM y la
        comparación dejaría de ser válida.
  \item Se limpia la caché y \textbf{se mide el efecto} (cuánto bajó
        \texttt{buff/cache}), no solo el código de salida del comando. Entre 0.7 y
        1.8~GB liberados por repetición, registrados en el archivo crudo.
  \item Se levanta el motor con el perfil de la celda y se espera su
        \emph{healthcheck}.
  \item \textbf{Se verifican los límites leyendo el cgroup} y se comparan con los
        esperados del perfil. Si no coinciden, la repetición se aborta sin
        escribir fila.
  \item Ventanas de \textbf{60~s de calentamiento, 300~s de medición y 60~s de
        enfriamiento}, como fija el manual. Los clientes trabajan en bucle cerrado,
        sin pausa entre consultas.
\end{enumerate}

La concurrencia 50 se reparte en \textbf{10 procesos de 5 hilos}, cada hilo con su
propia conexión. Se verificó que el motor ve los 50 clientes simultáneos
igual que con 50 procesos, y se eligió este reparto porque con 50 procesos el
runner ocupaba unos 2~GB y empujaba la VM a \emph{swap}, frente a unos 460~MB con
10 procesos. La decisión se apoya en el P50, el throughput, la concurrencia vista
por el motor, las conexiones y la memoria; \textbf{no} en el P95, que quedó
contaminado por el arranque en frío.

\subsection{Cómo se resumen las repeticiones (pregunta 5 del revisor)}
\label{sec:resumen}

\textbf{Tres repeticiones por celda}, y las tablas de este reporte dan la
\textbf{media} de las tres. Se eligió la media y no la mediana porque con tres
valores la mediana es simplemente el del medio y descarta los otros dos. Para que
el promedio no esconda dispersión, cada celda se acompaña de su coeficiente de
variación, y las celdas cuya dispersión supera el 15~\% o que muestran una
tendencia monótona se señalan una por una en la sección~\ref{sec:resultados}.

\subsection{Cómo se calculan los percentiles (pregunta 2 del revisor)}

Se reportan \textbf{P50, P95 y P99}, no promedios de latencia. Se calculan con
interpolación lineal (\texttt{statistics.quantiles} con
\texttt{method="inclusive"}, el mismo método que \texttt{numpy} usa por defecto)
sobre las latencias de las consultas \textbf{completadas dentro de la ventana de
medición}. Las latencias crudas de cada consulta se conservan en
\texttt{results/raw/}, así que los percentiles se pueden recalcular con otro
método sin volver a medir.

\paragraph{Los \emph{timeouts} son datos censurados, no ceros.}
Una consulta que agota el límite de 300~s no se descarta ni se cuenta como un
valor: es un dato \textbf{censurado por la derecha}, del que solo se sabe que su
latencia es mayor que el límite. Por eso, cuando la fracción de \emph{timeouts}
dentro de la ventana alcanza el 1~\%, la fila se marca con
\texttt{status = error} y \textbf{cada percentil se declara válido o no}: el P99
solo vale si los \emph{timeouts} son menos del 1~\%, el P95 menos del 5~\% y el
P50 menos del 50~\%. Esa validez va fila por fila en \texttt{results/raw/}. En
este experimento \cifraCeldasCensuradas{} celdas quedaron con percentiles
invalidados, y se marcan en todas las tablas.

\paragraph{Celdas con pocas muestras.}
Con 300~s de ventana y consultas de 20~s, una celda puede completar 15
consultas. Se marcan como \textbf{baja densidad muestral} (menos de 20
completadas) y su P99 se lee con esa reserva, en vez de presentarlo como si
viniera de cientos de muestras.

\subsection{De dónde sale cada columna del CSV}

Las 14 columnas se miden directamente, y \textbf{ninguna se rellena con un cero
inventado}: si una lectura falta, la celda queda vacía y la fila sale con
\texttt{status} \texttt{error} u \texttt{oom}.

\begin{table}[htbp]
  \centering
  \caption{Origen de las columnas que no son identificadores.}
  \label{tab:origen-columnas}
  \small
  \begin{tabularx}{\linewidth}{lX}
    \toprule
    \texttt{throughput\_qps} & consultas completadas en la ventana $\div$ 300~s \\
    \texttt{latency\_p50/p95/p99\_ms} & percentiles de las completadas en la ventana \\
    \texttt{cpu\_throttled\_seconds} & $\Delta$ \texttt{throttled\_usec} de
      \texttt{cpu.stat} del cgroup $\div$ \num{1e6} \\
    \texttt{max\_rss\_mb} & máximo de \texttt{anon} de \texttt{memory.stat},
      muestreado cada segundo durante la ventana \\
    \texttt{disk\_spill} & $\Delta$ \texttt{temp\_bytes} (PostgreSQL) o
      contadores de \emph{spill} de \texttt{serverStatus} (MongoDB) \\
    \texttt{status} & \texttt{oom} > \texttt{error} > \texttt{spill} > \texttt{ok} \\
    \bottomrule
  \end{tabularx}
\end{table}

El cgroup se lee \textbf{desde fuera del contenedor}, en
\texttt{/sys/fs/cgroup/docker/}, y no con \texttt{docker exec}: se midió que cada
lectura por \texttt{exec} cuesta unos 250~ms y gasta CPU \emph{del motor}, y
muestrear así cada segundo en el perfil C se comería un 8~\% de su cuota. Es
decir, mediríamos la medición.

\paragraph{Jerarquía de fuentes.}
Las columnas del CSV salen del cgroup y de los motores. Prometheus, cAdvisor y
node\_exporter (con \texttt{scrape\_interval: 1s}) alimentan las curvas y sirven
de control cruzado, pero \textbf{no} producen ninguna columna. Se verificó que
las dos fuentes miden lo mismo: \texttt{container\_memory\_rss} de cAdvisor es
exactamente \texttt{anon} del cgroup —comprobado con el contenedor congelado, donde
ningún proceso puede alterar la memoria— y el throttling de las dos fuentes
difiere menos del 1~\% bajo carga.

\subsection{Detección de suspensión del equipo}

Si el equipo se suspende durante una ventana, la VM se congela: las consultas en
vuelo se detienen y el reloj de pared avanza sin que corra nada. Esa ventana ya no
mide lo que dice medir. El runner lo detecta comparando el reloj de pared con el
monótono durante la ventana y marca la repetición como \texttt{error}. En esta
matriz \textbf{no ocurrió ninguna vez}, y el número de repeticiones descartadas
por cualquier motivo es \cifraDescartadas{} (sección~\ref{sec:integridad}).

% 03-resultados.tex — Sección 3: resultados consolidados.
% Responde las preguntas 2 y 3 del revisor. Las cifras vienen de cifras/cifras.tex
% y el cuerpo de la tabla de tablas/resultados.tex: nada escrito a mano.

\section{Resultados consolidados}
\label{sec:resultados}

\subsection{Qué se consiguió medir}

La matriz son \cifraCeldas{} celdas (2 motores $\times$ 3 perfiles $\times$ 2
consultas $\times$ 3 concurrencias) con 3 repeticiones cada una:
\cifraFilasCsv{} filas, todas en \texttt{results/measurements.csv}, validadas
contra el esquema exacto de 14 columnas del manual. Ninguna combinación falta ni
está repetida, y se descartaron \cifraDescartadas{} repeticiones.

\begin{table}[htbp]
  \centering
  \caption{Reparto de las \cifraFilasCsv{} filas por \texttt{status}. La
           precedencia es \texttt{oom} > \texttt{error} > \texttt{spill} >
           \texttt{ok}, así que una fila \texttt{oom} puede tener además volcado a
           disco: por eso \cifraFilasConSpill{} filas tienen
           \texttt{disk\_spill = true} y solo \cifraFilasSpill{} salen con
           \texttt{status = spill}.}
  \label{tab:estados}
  \small
  \begin{tabular}{lrl}
    \toprule
    \texttt{status} & Filas & Qué significa \\
    \midrule
    \texttt{ok}    & \cifraFilasOk{}    & la celda se midió sin incidencias \\
    \texttt{spill} & \cifraFilasSpill{} & hubo volcado a disco durante la ventana \\
    \texttt{error} & \cifraFilasError{} & percentiles invalidados por censura (A4) \\
    \texttt{oom}   & \cifraFilasOom{}   & el kernel mató algún proceso del motor \\
    \bottomrule
  \end{tabular}
\end{table}

De las \cifraCeldas{} celdas, \cifraCeldasMedibles{} son \textbf{medibles} (sus
tres repeticiones salieron \texttt{ok} o \texttt{spill}). Las otras cuatro
contienen filas \texttt{error} u \texttt{oom}: en ellas el motor no completó
trabajo suficiente como para que un percentil signifique algo, y
\textbf{eso es el resultado}, no un fallo del instrumento.

\subsection{Tabla consolidada}

\begin{longtable}{llllrrrrrrll}
  \caption{Resultados por celda: media de las 3 repeticiones. Latencias en ms,
           throttling en segundos por ventana de 300~s, \texttt{max\_rss\_mb} en
           MiB.\label{tab:resultados}}\\
  \toprule
  Motor & Perfil & Consulta & c & qps & P50 & P95 & P99 & Throttl. & RSS & Spill & \texttt{status} \\
  \midrule
  \endfirsthead
  \toprule
  Motor & Perfil & Consulta & c & qps & P50 & P95 & P99 & Throttl. & RSS & Spill & \texttt{status} \\
  \midrule
  \endhead
  \bottomrule
  \endfoot
  % generado por scripts/build_report_data.py — no editar a mano
MongoDB & A & Q1 & 1 & 0.2911 & 3414 & 3743 & 4011 & 0.0 & 823.6 & no & ok \\
MongoDB & A & Q1 & 10 & 0.5444 & 18149 & 21824 & 23459 & 1180.8 & 872.9 & no & ok \\
MongoDB & A & Q1 & 50 & 0.6856 & 74756 & 83773 & 85571 & 1185.2 & 1040.0 & no & ok \\
MongoDB & A & Q2 & 1 & 0.7300 & 1352 & 1526 & 1704 & 0.0 & 831.9 & no & ok \\
MongoDB & A & Q2 & 10 & 1.9078 & 5110 & 6979 & 7779 & 1186.5 & 861.2 & no & ok \\
MongoDB & A & Q2 & 50 & 2.2711 & 21775 & 24751 & 26324 & 1196.1 & 961.9 & no & ok \\
MongoDB & B & Q1 & 1 & 0.2611 & 3818 & 4019 & 4167 & 12.8 & 812.2 & no & ok \\
MongoDB & B & Q1 & 10 & 0.3300 & 30597 & 35030 & 37341 & 1783.6 & 822.0 & no & ok \\
MongoDB & B & Q1 & 50 & 0.3333 & 138931 & 148438 & 151839 & 1788.0 & 879.4 & no & ok \\
MongoDB & B & Q2 & 1 & 0.1800 & 5536 & 5934 & 6130 & 17.4 & 829.8 & no & ok \\
MongoDB & B & Q2 & 10 & 0.4111 & 23692 & 26105 & 26166 & 1502.9 & 805.5 & no & ok \\
MongoDB & B & Q2 & 50 & 0.3733 & 146673 & 172727 & 179510 & 1773.4 & 908.8 & no & ok \\
MongoDB & C & Q1 & 1 & 0.1933 & 5505 & 5956 & 6082 & 158.2 & 587.7 & no & ok \\
MongoDB & C & Q1 & 10 & 0.1800 & 57199 & 63125 & 64473 & 2081.3 & 607.3 & no & ok \\
MongoDB & C & Q1 & 50 & 0.1356 & 283597 & 296988 & 298965 & 2085.2 & 634.7 & no & error\tnotaCensurada{} \\
MongoDB & C & Q2 & 1 & 0.0533 & 19329 & 20461 & 20683 & 183.4 & 611.6 & no & ok\tnotaBajaDensidad{} \\
MongoDB & C & Q2 & 10 & 0.1000 & 96905 & 113894 & 113895 & 1830.0 & 561.6 & no & ok \\
MongoDB & C & Q2 & 50 & 0.1667 & 225253 & 225372 & 228109 & 1976.7 & 672.3 & no & ok \\
PostgreSQL & A & Q1 & 1 & 0.1322 & 7481 & 8321 & 9030 & 0.3 & 103.8 & sí & spill \\
PostgreSQL & A & Q1 & 10 & 0.1444 & 74687 & 97188 & 98180 & 826.9 & 570.9 & sí & spill \\
PostgreSQL & A & Q1 & 50 & 0.0022 & 275776 & 283670 & 284371 & 552.6 & 1851.6 & sí & error\tnotaCensurada{} \\
PostgreSQL & A & Q2 & 1 & 1.8800 & 513 & 715 & 795 & 18.4 & 15.6 & no & ok \\
PostgreSQL & A & Q2 & 10 & 2.3511 & 3507 & 7477 & 8168 & 907.7 & 43.4 & no & ok \\
PostgreSQL & A & Q2 & 50 & 1.5789 & 35086 & 43086 & 50592 & 1093.6 & 127.3 & no & ok \\
PostgreSQL & B & Q1 & 1 & 0.1100 & 8923 & 10147 & 11198 & 146.0 & 105.0 & sí & spill \\
PostgreSQL & B & Q1 & 10 & 0.1000 & 114076 & 118444 & 119141 & 1535.2 & 571.0 & sí & spill \\
PostgreSQL & B & Q1 & 50 & 0.0133 & 208046 & 262901 & 270597 & 1216.0 & 1800.0 & sí & error/oom\tnotaCensurada{} \\
PostgreSQL & B & Q2 & 1 & 1.0889 & 885 & 1130 & 1311 & 550.0 & 16.3 & no & ok \\
PostgreSQL & B & Q2 & 10 & 1.1478 & 7330 & 15423 & 17939 & 1379.2 & 43.4 & no & ok \\
PostgreSQL & B & Q2 & 50 & 1.0556 & 51582 & 59091 & 62590 & 1634.3 & 127.3 & no & ok \\
PostgreSQL & C & Q1 & 1 & 0.0567 & 17555 & 18464 & 18777 & 489.3 & 104.4 & sí & spill\tnotaBajaDensidad{} \\
PostgreSQL & C & Q1 & 10 & 0.0333 & 263918 & 267051 & 267746 & 1829.8 & 564.7 & sí & spill\tnotaBajaDensidad{} \\
PostgreSQL & C & Q1 & 50 & 0.0000 & --- & --- & --- & 105.5 & 475.1 & no & oom\tnotaColapsada{} \\
PostgreSQL & C & Q2 & 1 & 0.5722 & 1730 & 2024 & 2108 & 870.7 & 16.7 & no & ok \\
PostgreSQL & C & Q2 & 10 & 0.6000 & 11936 & 28847 & 29888 & 1622.2 & 43.8 & no & ok \\
PostgreSQL & C & Q2 & 50 & 0.5611 & 100284 & 124632 & 134570 & 1890.9 & 127.0 & no & ok \\%

\end{longtable}

\leyendaNotas{}

\subsection{Throttling de CPU (pregunta 3 del revisor)}

Sí hubo throttling, y es la métrica que más fácil se malinterpreta. Dos
advertencias necesarias para leer la columna:

\begin{enumerate}
  \item \textbf{Puede superar la duración de la ventana.} El kernel suma el tiempo
        de throttling de \textbf{cada CPU} en la que el cgroup tenía hilos listos
        para correr. Con 4 CPUs y una cuota de 1, cada segundo de reloj puede
        acumular hasta 4 segundos de throttling. El máximo observado es de
        \textbf{2086~s en una ventana de 300~s}. No es un error de medida: cAdvisor
        reporta la misma suma, y las dos fuentes coinciden al 1~\%.
  \item \textbf{Una parte no la causa el benchmark.} Se midió el throttling de cada
        motor \textbf{en reposo}, sin consultas, por perfil: MongoDB acumula unos
        27~s por ventana de 300~s en el perfil~C y 3.6~s en el~B, mientras
        PostgreSQL se queda por debajo de 0.5~s en todos. Esa línea base está en
        \texttt{results/f6/f6\_3/baseline/} y aparece en la
        figura~\ref{fig:throttling-qps}.
\end{enumerate}

\subsection{Estabilidad de las mediciones}
\label{sec:estabilidad}

Antes de gastar las horas de la matriz se corrió una celda piloto completa y se
exigió que el coeficiente de variación del throughput entre las tres repeticiones
fuera menor del 15~\%; salió del 4.2~\% en MongoDB y del 3.8~\% en PostgreSQL.
Sobre la matriz entera, \textbf{la dispersión se mantuvo baja en casi todas las
celdas}, pero no en todas, y conviene decir cuáles:

\begin{itemize}
  \item \textbf{Dos celdas superan el 15~\%}, las dos de \texttt{Q2} con
        \mbox{c=50}: PostgreSQL en el perfil~A (20.9~\%, con una caída del 34~\%
        entre la primera y la tercera repetición) y MongoDB en el perfil~B
        (24.0~\%, con una caída del 33~\%).
  \item En la de \textbf{PostgreSQL} la causa es medible: las tres repeticiones
        gastaron \emph{el mismo} tiempo de CPU (1197~s) y el rendimiento del
        procesador cayó del 75.9~\% al 65.2~\% justo en la repetición que pierde el
        34~\%. El mismo CPU, más lento, hace menos trabajo.
  \item En la de \textbf{MongoDB} la causa \textbf{no está identificada}: el
        tiempo de CPU es idéntico (598.6~s) y el reloj está plano (83.4~\%,
        81.7~\%, 84.6~\%). La caché de Windows sobre el disco virtual, que
        \texttt{drop\_caches} no alcanza, es el sospechoso natural, pero no hay
        medición que lo demuestre y no se presenta como explicación.
  \item Se comprobó que \textbf{no hay un sesgo sistemático} de «la primera
        repetición es más rápida»: ocurre en 6 de las 30 celdas con las tres
        repeticiones comparables, cuando por azar serían unas 10.
\end{itemize}

La consecuencia práctica es que la celda de \texttt{Q2} con \mbox{c=50} en el
perfil~A es la \textbf{única} en la que MongoDB supera a PostgreSQL en
\texttt{Q2}, y es también una de las dos celdas con dispersión alta. Esa
comparación concreta no se usa para sostener ninguna conclusión
(sección~\ref{sec:interpretacion}).

% 04-graficas.tex — Sección 4: análisis gráfico (R-30). Las tres figuras
% obligatorias, dibujadas por pgfplots desde report/datos/*.dat.

\section{Análisis gráfico}
\label{sec:graficas}

Las tres figuras las dibuja LaTeX leyendo los archivos de
\texttt{report/datos/}, que genera \texttt{scripts/build\_report\_data.py} del
CSV y de \texttt{results/raw/}. Así cada punto es rastreable hasta su fila.

% p99-concurrencia.tex — Figura obligatoria 1 (R-30): P99 contra concurrencia.
%
% R-30.1 (decisión de Alix, al aprobar F7-e): la figura muestra LAS DOS VISTAS, no
% una nota al pie. Con c=50 en los perfiles B y C, hasta el 91 % de las consultas
% de la ventana son la primera consulta en frío de cada cliente (F7.3), así que
% parte de la subida de la curva es arranque en frío y no latencia del motor.
%   - línea continua: el P99 de la columna del CSV (lo que mide el manual);
%   - línea punteada: el P99 recalculado sin la primera consulta de cada cliente,
%     desde results/raw/, que es la vista en régimen.
% Donde el .dat trae `nan` no hay vista en régimen porque TODAS las consultas de
% la ventana eran primeras consultas: la línea punteada se interrumpe ahí
% (unbounded coords=jump, en el preámbulo).
%
% Los cuatro paneles van escritos uno a uno, SIN un macro con parámetros: dentro
% de un \newcommand, los '#' del cuerpo chocan con el analizador de claves de
% pgfplots ("Illegal parameter number in definition of \pgfkeys@..."). Es más
% largo, pero no depende de cómo pgfkeys reinterprete el cuerpo.

\begin{figure}[htbp]
  \centering
  \begin{tikzpicture}
    \begin{axis}[
        title={MongoDB · Q2 · perfil A}, title style={font=\small},
        xlabel={Concurrencia (clientes)}, ylabel={P99 (ms)},
        xmode=log, log basis x=10, xtick={1,10,50}, xticklabels={1,10,50},
        ymin=0, width=0.47\linewidth, height=5cm,
        legend pos=north west, legend cell align=left,
      ]
      \addplot[color=blue!60!black, mark=*, thick] table [x=concurrencia, y=p99_csv] {datos/p99_mongodb_Q2_A.dat};
      \addplot[color=blue!60!black, mark=o, densely dotted, thick]
        table [x=concurrencia, y=p99_sin_frias] {datos/p99_mongodb_Q2_A.dat};
      \legend{P99 del CSV, Sin la 1.\textsuperscript{a} en frío}
    \end{axis}
  \end{tikzpicture}\hfill
  \begin{tikzpicture}
    \begin{axis}[
        title={PostgreSQL · Q2 · perfil A}, title style={font=\small},
        xlabel={Concurrencia (clientes)}, ylabel={P99 (ms)},
        xmode=log, log basis x=10, xtick={1,10,50}, xticklabels={1,10,50},
        ymin=0, width=0.47\linewidth, height=5cm,
        legend pos=north west, legend cell align=left,
      ]
      \addplot[color=red!60!black, mark=*, thick] table [x=concurrencia, y=p99_csv] {datos/p99_postgresql_Q2_A.dat};
      \addplot[color=red!60!black, mark=o, densely dotted, thick]
        table [x=concurrencia, y=p99_sin_frias] {datos/p99_postgresql_Q2_A.dat};
    \end{axis}
  \end{tikzpicture}

  \begin{tikzpicture}
    \begin{axis}[
        title={MongoDB · Q2 · perfil C}, title style={font=\small},
        xlabel={Concurrencia (clientes)}, ylabel={P99 (ms)},
        xmode=log, log basis x=10, xtick={1,10,50}, xticklabels={1,10,50},
        ymin=0, width=0.47\linewidth, height=5cm,
        legend pos=north west, legend cell align=left,
      ]
      \addplot[color=blue!60!black, mark=*, thick] table [x=concurrencia, y=p99_csv] {datos/p99_mongodb_Q2_C.dat};
      \addplot[color=blue!60!black, mark=o, densely dotted, thick]
        table [x=concurrencia, y=p99_sin_frias] {datos/p99_mongodb_Q2_C.dat};
    \end{axis}
  \end{tikzpicture}\hfill
  \begin{tikzpicture}
    \begin{axis}[
        title={PostgreSQL · Q2 · perfil C}, title style={font=\small},
        xlabel={Concurrencia (clientes)}, ylabel={P99 (ms)},
        xmode=log, log basis x=10, xtick={1,10,50}, xticklabels={1,10,50},
        ymin=0, width=0.47\linewidth, height=5cm,
        legend pos=north west, legend cell align=left,
      ]
      \addplot[color=red!60!black, mark=*, thick] table [x=concurrencia, y=p99_csv] {datos/p99_postgresql_Q2_C.dat};
      \addplot[color=red!60!black, mark=o, densely dotted, thick]
        table [x=concurrencia, y=p99_sin_frias] {datos/p99_postgresql_Q2_C.dat};
    \end{axis}
  \end{tikzpicture}

  \caption[P99 contra concurrencia, en dos vistas]{%
    \textbf{P99 contra concurrencia, en dos vistas (R-30.1).} La línea continua es
    el P99 tal como lo mide el manual (columna del CSV). La punteada es el mismo
    P99 \textbf{sin la primera consulta de cada cliente}, recalculado desde
    \texttt{results/raw/}. La separación entre las dos líneas es el arranque en
    frío: el calentamiento de 60~s del manual no alcanza con \mbox{c=50} en los
    perfiles B y C, donde hasta el 91~\% de la ventana son primeras consultas
    (F7.3). En MongoDB · perfil C · \mbox{c=50} la línea punteada se interrumpe
    porque \emph{todas} las consultas de la ventana eran primeras consultas y no
    existe un P99 en régimen que dibujar. \textbf{Leer solo la línea continua
    sobreestima la latencia del motor en esos puntos.} Datos:
    \texttt{report/datos/p99\_*.dat}.}
  \label{fig:p99-concurrencia}
\end{figure}


La figura~\ref{fig:p99-concurrencia} lleva \textbf{dos series por panel} y no una
nota al pie, porque la diferencia entre ellas no es un detalle: el calentamiento
de 60~s que fija el manual no alcanza cuando 50 clientes empiezan a la vez con la
caché recién limpiada. Se midió una por una la primera consulta de cada cliente
y, con \mbox{c=50}, \textbf{no cabe en los 60~s en 14 de las 18 combinaciones};
solo cabe en el perfil~A con \texttt{Q2} (de 39 a 52~s). En los perfiles B y C
llega a 110~s, y en algunas celdas ninguna consulta terminó siquiera en los 110~s
que dura toda la corrida de prueba.

El caso extremo aparece en la figura: en \textbf{MongoDB, perfil~C,
\texttt{Q2}, \mbox{c=50}} las \num{150} consultas que terminaron dentro de la
ventana son \emph{todas} la primera de su cliente, así que no existe un P99 «en
régimen» que dibujar y la serie punteada se interrumpe. Decir que el P99 de esa
celda es de 228~s sin esta advertencia sería atribuirle al motor el costo de
arrancar en frío.

% throttling-qps.tex — Figura obligatoria 2 (R-30): throttling contra throughput.
%
% Lo que hay que saber para leer el eje (medido en F6.2, declarado en SPEC §4):
% `cpu_throttled_seconds` es la SUMA por CPU del tiempo que el cgroup tuvo hilos
% listos y no pudo correr. Por eso puede superar la duración de la ventana: con 4
% CPUs y una cuota de 1, cada segundo de reloj acumula hasta 4 de throttling. El
% eje NO se acota a 300 s.
%
% Además (F6.3): parte del throttling de MongoDB en el perfil C ocurre en reposo,
% sin benchmark (~27 s por ventana de 300 s). La línea de referencia lo marca.
%
% Un archivo de datos por motor y concurrencia (report/datos/throttling_*.dat),
% para no depender de filtros de pgfplots.

\begin{figure}[htbp]
  \centering
  \begin{tikzpicture}
    \begin{axis}[
        xlabel={Throttling de CPU en la ventana de 300~s (s)},
        ylabel={Throughput (consultas/s)},
        ymode=log, log basis y=10,
        xmin=0,   % el throttling no puede ser negativo
        enlarge x limits={upper, value=0.05},
        legend pos=south west, legend cell align=left, legend columns=2,
        height=7.5cm, width=\linewidth,
      ]
      \addplot[only marks, mark=*,       mark size=1.6pt, blue!65!black]
        table [x=throttling_s, y=qps] {datos/throttling_mongodb_c1.dat};
      \addplot[only marks, mark=*,       mark size=2.6pt, blue!65!black]
        table [x=throttling_s, y=qps] {datos/throttling_mongodb_c10.dat};
      \addplot[only marks, mark=*,       mark size=3.6pt, blue!65!black]
        table [x=throttling_s, y=qps] {datos/throttling_mongodb_c50.dat};
      \addplot[only marks, mark=square*, mark size=1.6pt, red!65!black]
        table [x=throttling_s, y=qps] {datos/throttling_postgresql_c1.dat};
      \addplot[only marks, mark=square*, mark size=2.6pt, red!65!black]
        table [x=throttling_s, y=qps] {datos/throttling_postgresql_c10.dat};
      \addplot[only marks, mark=square*, mark size=3.6pt, red!65!black]
        table [x=throttling_s, y=qps] {datos/throttling_postgresql_c50.dat};
      \legend{MongoDB c=1, MongoDB c=10, MongoDB c=50,
              PostgreSQL c=1, PostgreSQL c=10, PostgreSQL c=50}
      % F6.3: línea base de throttling en reposo de MongoDB en el perfil C.
      \draw[dashed, gray!70] ({axis cs:26.7,0}|-{rel axis cs:0,0})
                          -- ({axis cs:26.7,0}|-{rel axis cs:0,1});
      \node[gray!70, font=\tiny, anchor=south west, rotate=90]
        at ({axis cs:26.7,0}|-{rel axis cs:0,0.05}) {MongoDB · C en reposo (F6.3)};
    \end{axis}
  \end{tikzpicture}
  \caption[Throttling contra throughput]{%
    \textbf{Throttling de CPU contra throughput.} Cada punto es una celda (media de
    sus 3 repeticiones); el tamaño de la marca crece con la concurrencia.
    \textbf{El eje horizontal supera los 300~s de la ventana} porque
    \texttt{cpu\_throttled\_seconds} suma el tiempo de cada CPU en la que el cgroup
    tenía hilos listos (medido en F6.2, declarado en SPEC §4); el máximo observado
    es de \num{2086}~s. La línea vertical marca el throttling que MongoDB acumula
    \textbf{en reposo} en el perfil C, sin benchmark (F6.3): a su izquierda no hay
    contención atribuible a la carga. Datos:
    \texttt{report/datos/throttling\_*.dat}.}
  \label{fig:throttling-qps}
\end{figure}


% curva-rss.tex — Figura obligatoria 3 (R-30): curva de memoria durante la ventana.
%
% Dos series por panel, y la diferencia entre ellas es el hallazgo F6-A3:
%   anon        = lo que mide la columna `max_rss_mb` del CSV;
%   working_set = memory.current − inactive_file, que SÍ incluye la memoria
%                 compartida (`shared_buffers`) de PostgreSQL.
% Los `shared_buffers` son `shmem`, no `anon`, así que con `max_rss_mb` PostgreSQL
% sale artificialmente eficiente (SPEC §4 y A11). M3 se interpreta con working_set.
%
% Fuente: la serie de 1 s del cgroup que el runner guarda en results/raw/
% (serie_cgroup_1s), es decir las MISMAS lecturas de las que sale la columna.

\newcommand{\rssPanel}[4]{%
  % #1 archivo .dat · #2 título · #3 color · #4 "leyenda" para mostrarla
  \begin{tikzpicture}
    \begin{axis}[
        title={#2}, title style={font=\small},
        xlabel={Segundo de la ventana de medición}, ylabel={Memoria (MiB)},
        xmin=0, xmax=300, ymin=0,
        width=0.47\linewidth, height=5cm,
        legend pos=south east, legend cell align=left,
      ]
      \addplot[color=#3, thick, no marks] table [x=t_s, y=working_set_mib] {#1};
      \addplot[color=#3!55, thick, densely dashed, no marks]
        table [x=t_s, y=anon_mib] {#1};
      \addplot[color=gray!60, thin, dotted, no marks] table [x=t_s, y=shmem_mib] {#1};
      \ifthenelse{\equal{#4}{leyenda}}%
        {\legend{\texttt{working\_set}, \texttt{anon}, \texttt{shmem}}}{}%
    \end{axis}
  \end{tikzpicture}%
}

\begin{figure}[htbp]
  \centering
  \rssPanel{rss_mongodb_A_Q2_c50.dat}{MongoDB · perfil A · Q2 · c=50}{blue!60!black}{leyenda}\hfill
  \rssPanel{rss_postgresql_A_Q2_c50.dat}{PostgreSQL · perfil A · Q2 · c=50}{red!60!black}{}

  \rssPanel{rss_mongodb_C_Q2_c50.dat}{MongoDB · perfil C · Q2 · c=50}{blue!60!black}{}\hfill
  \rssPanel{rss_postgresql_C_Q2_c50.dat}{PostgreSQL · perfil C · Q2 · c=50}{red!60!black}{}

  \caption[Curva de memoria durante la ventana]{%
    \textbf{Memoria del contenedor durante la ventana, muestreada cada segundo.}
    La línea discontinua (\texttt{anon}) es exactamente lo que resume la columna
    \texttt{max\_rss\_mb}; la continua (\texttt{working\_set}) incluye además la
    memoria compartida. \textbf{En PostgreSQL la distancia entre las dos es
    \texttt{shared\_buffers}}: son \texttt{shmem}, no \texttt{anon}, así que
    \texttt{max\_rss\_mb} no cuenta su búfer principal y lo hace parecer
    artificialmente eficiente frente a MongoDB, cuyo caché de WiredTiger sí es
    \texttt{anon} (F6-A3, SPEC §4 y A11). Por eso M3 se interpreta con
    \texttt{working\_set}. Fuente: serie de 1 s del cgroup en
    \texttt{results/raw/}, las mismas lecturas de las que sale la columna.}
  \label{fig:curva-rss}
\end{figure}


La figura~\ref{fig:curva-rss} es la que obliga a tener cuidado con la columna
\texttt{max\_rss\_mb}, y por eso lleva dos series: véase la
sección~\ref{sec:memoria}.

% 05-planes.tex — Sección 5: planes de ejecución. Responde la pregunta 6 del
% revisor y declara la asimetría de planes (riesgo R14 del PLAN).

\section{Planes de ejecución}
\label{sec:planes}

Se capturó el plan de cada combinación de motor, consulta y perfil con
\texttt{explain} en modo \texttt{executionStats} (MongoDB) y
\texttt{EXPLAIN (ANALYZE, BUFFERS)} (PostgreSQL). Los 12 planes crudos están en
\texttt{results/explains/}.

\subsection{Dos avisos para leerlos}

\paragraph{Los planes se capturan con un solo cliente.}
Se toman en las celdas de \mbox{c=1} y \textbf{fuera} de la ventana de medición,
para no alterar lo que se mide. Eso tiene una consecuencia que hay que declarar:
\textbf{el plan de \mbox{c=1} no es el que se ejecuta con \mbox{c=50}}. Con 50
clientes simultáneos, PostgreSQL agota los procesos de trabajo paralelos
disponibles y los clientes que llegan después ejecutan la consulta sin
paralelismo. El plan documenta la \emph{forma} de resolver la consulta, no el
grado de paralelismo que tuvo cada cliente bajo carga.

\paragraph{En \texttt{EXPLAIN (BUFFERS)} los contadores incluyen a los hijos.}
El total es el del nodo raíz. Sumar los nodos cuenta lo mismo varias veces: al
verificar las consultas se registraron 8310 bloques donde en realidad eran 1662.

\subsection{La asimetría de planes, declarada}
\label{sec:asimetria-planes}

\texttt{Q2} filtra por \texttt{timestamp} y ordena por \texttt{final\_total}. Los
dos motores la resuelven de forma distinta:

\begin{itemize}
  \item \textbf{MongoDB} usa su índice sobre \texttt{timestamp} para el
        \texttt{\$match}.
  \item \textbf{PostgreSQL} hace \texttt{Seq Scan}. El índice de expresión existe
        —se creó sobre el \texttt{timestamp} extraído como texto, sin modificar la
        consulta del manual— pero el planificador no lo usa, porque la mitad de la
        tabla pasa el filtro y recorrerla entera le sale más barato que saltar por
        el índice.
\end{itemize}

Esto es importante al interpretar las latencias: \textbf{en parte se están
comparando dos planes distintos}, no solo dos motores. No se «arregló» forzando
el índice, porque eso significaría medir una consulta que no es la del manual.

\subsection{Lo que el plan explica del cuello de botella (pregunta 6)}

\paragraph{\texttt{Q1} en PostgreSQL: el plan explica el volcado.}
El plan muestra el ordenamiento con \texttt{Sort Method: external merge} y
\texttt{Sort Space Type: Disk}, que es exactamente lo que registran los
contadores \texttt{temp\_bytes} de la sección~\ref{sec:spill}. Las dos fuentes
coinciden: no es una conjetura sobre por qué \texttt{Q1} es lenta en PostgreSQL.

\paragraph{\texttt{Q2} en los dos motores: el plan explica por qué \emph{no} hay volcado.}
El \texttt{LIMIT 1000} convierte el ordenamiento en un \emph{top-K}: no hace falta
ordenar las \num{175049} filas que pasan el filtro, solo mantener las 1000
mejores. PostgreSQL lo resuelve con \texttt{top-N heapsort} y
\texttt{Sort Space Used: 273}~kB \textbf{en memoria}. Esa cifra es la clave de la
sección~\ref{sec:variantes}: 273~kB caben de sobra tanto en los 4~MB de
\texttt{work\_mem} de la variante restrictiva como en los 64~MB de la permisiva,
y por eso las dos variantes de \texttt{Q2} no se diferencian.

% 06-interpretacion.tex — Sección 6: interpretación. Responde M1, M2 y M3 con
% datos, y aplica las tres decisiones de Alix: R-29.1 (distinguir dónde ocurrió
% cada OOM), R-29.2 (nota al pie de la celda colapsada) y R-16.1 (la premisa del
% manual sobre las variantes de Q2, con la cobertura parcial declarada).

\section{Interpretación}
\label{sec:interpretacion}

\subsection{M1: ¿cómo se degrada cada motor al apretar los recursos?}
\label{sec:m1}

El resultado no se parte por motor, se parte \textbf{por consulta}, y es
consistente en los tres perfiles y en las tres concurrencias.

\begin{table}[htbp]
  \centering
  \caption{Quién gana cada celda, en throughput (media de 3 repeticiones). El
           factor es cuántas veces más consultas por segundo completa el ganador.}
  \label{tab:m1}
  \small
  \begin{tabular}{llrrr}
    \toprule
    Consulta & Perfil & c=1 & c=10 & c=50 \\
    \midrule
    \texttt{Q1} & A & Mongo $\times$2.2 & Mongo $\times$3.8 & PostgreSQL no completa\tnotaCensurada{} \\
    \texttt{Q1} & B & Mongo $\times$2.4 & Mongo $\times$3.3 & Mongo $\times$25 \\
    \texttt{Q1} & C & Mongo $\times$3.4 & Mongo $\times$5.4 & los dos fallan\tnotaColapsada{} \\
    \texttt{Q2} & A & PostgreSQL $\times$2.6 & PostgreSQL $\times$1.2 & Mongo $\times$1.4\tnotaBajaDensidad{} \\
    \texttt{Q2} & B & PostgreSQL $\times$6.0 & PostgreSQL $\times$2.8 & PostgreSQL $\times$2.8 \\
    \texttt{Q2} & C & PostgreSQL $\times$10.7 & PostgreSQL $\times$6.0 & PostgreSQL $\times$3.4 \\
    \bottomrule
  \end{tabular}
\end{table}

\paragraph{\texttt{Q1} (desanidar y agrupar): MongoDB gana siempre.}
Y la ventaja \textbf{crece} al apretar los recursos: de $\times$2.2 en el perfil A
a $\times$3.4 en el C con un cliente. La razón está en el plan y en el disco: para
PostgreSQL, \texttt{Q1} exige un ordenamiento que no cabe en \texttt{work\_mem} y
termina en disco \textbf{en los tres perfiles} (sección~\ref{sec:spill}), mientras
MongoDB la resuelve sin tocar disco.

\paragraph{\texttt{Q2} (ordenar y quedarse con el Top-1000): PostgreSQL gana, y la ventaja crece al apretar.}
De $\times$2.6 a $\times$10.7 con un cliente al pasar del perfil A al C. Visto al
revés: entre el perfil A y el C, MongoDB se vuelve \textbf{13.7 veces} más lento
en \texttt{Q2} y solo 1.5 veces más lento en \texttt{Q1}. La sensibilidad a la
contención depende más de la consulta que del motor.

\paragraph{La única excepción, y por qué no se usa.}
La celda \texttt{Q2} · perfil~A · \mbox{c=50} es la única de \texttt{Q2} que gana
MongoDB, y es también \textbf{una de las dos celdas con dispersión por encima del
15~\%} (sección~\ref{sec:estabilidad}): su medición de PostgreSQL cayó un 34~\%
entre la primera y la tercera repetición mientras el procesador bajaba de
frecuencia. Esa inversión no se usa para sostener nada.

\subsection{M2: ¿hubo volcado a disco y OOM, y dónde?}
\label{sec:spill}

\paragraph{Volcado: sí, masivo, y solo en PostgreSQL con \texttt{Q1}.}
\cifraFilasConSpill{} de las \cifraFilasCsv{} filas tienen
\texttt{disk\_spill = true}, y \textbf{todas son de PostgreSQL con \texttt{Q1}}:
8 de sus 9 celdas (la novena no volcó porque el motor murió antes). En total,
\textbf{\cifraTempTotalTB{}~TB} de archivos temporales escritos a lo largo de la
matriz. Con un solo cliente en el perfil A son unos 32~GB por ventana de 300~s,
es decir \textbf{unos 0.8~GB por consulta}.

\textbf{MongoDB no volcó en ninguna celda}, ni siquiera con
\texttt{allowDiskUse: true} (anexo~\ref{sec:variantes}).

\paragraph{OOM: cinco filas, y hay que distinguir dónde ocurrió cada uno.}
Son \cifraFilasOom{} filas con \texttt{status = oom}, todas de PostgreSQL con
\texttt{Q1} y \mbox{c=50} en los perfiles B y C. Presentarlas como «cinco OOM» sin
más sería impreciso:

\begin{itemize}
  \item \textbf{\cifraOomEnVentana{} ocurrieron dentro de la ventana de medición}:
        el contador \texttt{oom\_kill} del cgroup se movió mientras se medía.
  \item \textbf{\cifraOomEnCalentamiento{} ocurrieron durante el calentamiento}: el
        contador ya estaba en su valor final al abrirse la ventana. El motor llegó
        a la medición \emph{ya} sin memoria.
\end{itemize}

La distinción importa porque cambia lo que la fila mide. Las tres de la segunda
clase no describen «un motor que se queda sin memoria atendiendo la carga medida»,
sino «un motor que ya se había quedado sin memoria antes de empezar a medir». El
detalle de cada una está en \texttt{results/logs/}.

\begin{table}[htbp]
  \centering
  \caption{Las \cifraFilasOom{} repeticiones con \texttt{status = oom}. La columna
           «$\Delta$ \texttt{oom\_kill}» es cuántos procesos mató el kernel
           \emph{durante la ventana}: un cero con \texttt{OOMKilled = sí}
           significa que el OOM ocurrió antes, en el calentamiento (R-29.1).}
  \label{tab:oom}
  \small
  \begin{tabular}{lllrrrlrl}
    \toprule
    Motor & Perfil & Consulta & c & Rep. & $\Delta$ \texttt{oom\_kill} & Cuándo &
    Completadas & \texttt{OOMKilled} \\
    \midrule
    % generado por scripts/build_report_data.py — no editar a mano
PostgreSQL & B & Q1 & 50 & 1 & 0 & en el calentamiento & 4 & sí \\
PostgreSQL & B & Q1 & 50 & 2 & 0 & en el calentamiento & 4 & sí \\
PostgreSQL & C & Q1 & 50 & 1 & 0 & en el calentamiento & 0 & sí \\
PostgreSQL & C & Q1 & 50 & 2 & 1 & en la ventana & 0 & sí \\
PostgreSQL & C & Q1 & 50 & 3 & 6 & en la ventana & 0 & sí \\
   % el archivo generado cierra con \bottomrule (ver report/README.md)
  \end{tabular}
\end{table}

\paragraph{La peor celda, y cómo no leerla.}
En \textbf{PostgreSQL · perfil C · \texttt{Q1} · \mbox{c=50}} no se completó
\emph{ninguna} consulta en las tres repeticiones: el throughput es 0.000. La fila
registra \texttt{cpu\_throttled\_seconds = 0.035} y
\texttt{max\_rss\_mb = 8.2}\tnotaColapsada{}, y esos dos números \textbf{no
significan que PostgreSQL atendiera 50 clientes con 8~MiB de memoria}: significan
que durante la ventana el motor estaba colapsado y no trabajaba. Sus
\num{15000} errores registrados son \textbf{conexiones rechazadas}, no consultas
ejecutadas. Por eso esa celda lleva nota al pie en todas las tablas de este
reporte.

\paragraph{Respuesta a M2.} El volcado a disco es real, sistemático y enorme, pero
está concentrado en una sola combinación: PostgreSQL con \texttt{Q1}. El OOM
aparece solo donde se juntan la consulta más pesada, la concurrencia más alta y
los perfiles con menos memoria. MongoDB no volcó ni fue matado por el kernel en
ninguna celda.

\subsection{La premisa del manual sobre las dos variantes de \texttt{Q2}}
\label{sec:variantes}

El manual pide dos variantes de estrés de memoria para \texttt{Q2}
(\texttt{work\_mem} de 4~MB frente a 64~MB en PostgreSQL;
\texttt{allowDiskUse} falso frente a verdadero en MongoDB) \textbf{dando por
supuesto que se diferenciarían}. Se midieron las dos en el perfil con menos
memoria y \textbf{no se diferencian}: ni en throughput más allá de la dispersión
propia de cada celda, ni en volcado, que no aparece en ninguna de las dos.

La explicación está en el plan (sección~\ref{sec:planes}): el
\texttt{LIMIT 1000} convierte el ordenamiento en un \emph{top-K} de
\textbf{273~kB}, que caben de sobra en los 4~MB de la variante restrictiva. El
parámetro que el manual varía nunca llega a ser el límite. El detalle de la
comparación está en el anexo~\ref{sec:variantes-anexo}, con la cobertura parcial
declarada.

\subsection{M3: huella de CPU y memoria por cada 1000 documentos}
\label{sec:memoria}

El CSV no tiene un denominador, así que se toma de los \textbf{planes de
ejecución}, como corresponde: los documentos que cada consulta examina de verdad.
Y ahí aparece una asimetría que hay que declarar antes de dividir:

\begin{center}
  \small
  \begin{tabular}{llr}
    \toprule
    Motor & Consulta & Documentos examinados por consulta \\
    \midrule
    MongoDB    & \texttt{Q1} & \num{500000} (recorrido completo) \\
    MongoDB    & \texttt{Q2} & \num{250207} (usa su índice de \texttt{timestamp}) \\
    PostgreSQL & \texttt{Q1} & \num{500000} \\
    PostgreSQL & \texttt{Q2} & \num{500000} (\texttt{Seq Scan}: no usa el índice) \\
    \bottomrule
  \end{tabular}
\end{center}

\paragraph{CPU por 1000 documentos (perfil A, un cliente).}
MongoDB gasta 8.6~ms en \texttt{Q1} y PostgreSQL 25.6~ms: \textbf{MongoDB es unas
tres veces más eficiente en CPU} para esa consulta. En \texttt{Q2} se invierte:
PostgreSQL gasta 3.5~ms y MongoDB 6.6~ms. Y al apretar los recursos la brecha de
\texttt{Q2} se abre hasta ser de \textbf{74.7~ms frente a 3.6~ms} en el perfil~C,
unas 21 veces.

\paragraph{La memoria depende de cuál de las dos memorias se mire, y esto es importante.}
\texttt{max\_rss\_mb} mide la memoria \emph{anónima} del contenedor. Los
\texttt{shared\_buffers} de PostgreSQL son memoria \textbf{compartida}, que el
kernel no cuenta como anónima; el caché de WiredTiger de MongoDB sí lo es. El
efecto es grande:

\begin{center}
  \small
  \begin{tabular}{lrrrr}
    \toprule
    & \multicolumn{2}{c}{\texttt{max\_rss\_mb} (anónima)} & \multicolumn{2}{c}{\texttt{working\_set}} \\
    \cmidrule(lr){2-3}\cmidrule(lr){4-5}
    Perfil A, c=1 & \texttt{Q1} & \texttt{Q2} & \texttt{Q1} & \texttt{Q2} \\
    \midrule
    MongoDB    & 824 & 832 & 1042 & 1058 \\
    PostgreSQL & 104 & 16  & 1630 & 891 \\
    \bottomrule
  \end{tabular}
\end{center}

Leyendo solo \texttt{max\_rss\_mb}, PostgreSQL parecería de 8 a 50 veces más
eficiente en memoria. Leyendo \texttt{working\_set} —que sí incluye la memoria
compartida y es la que cuenta contra el límite del contenedor— \textbf{la relación
se invierte en \texttt{Q1}}: 1630~MiB de PostgreSQL frente a 1042 de MongoDB. Por
eso M3 se responde con \texttt{working\_set}, y la figura~\ref{fig:curva-rss}
muestra las dos curvas juntas.

\paragraph{Un matiz honesto sobre \texttt{working\_set}.}
Incluye también la caché de páginas de los archivos de datos que se le atribuye al
contenedor. PostgreSQL depende de ella por diseño —lee a través de la caché del
sistema operativo— y cuenta contra su \texttt{memory.max}, así que es memoria que
el motor necesita bajo el límite; pero no es «memoria del proceso». Las dos
lecturas están en los datos (\texttt{report/datos/m3\_huella.dat}) para que
cualquiera pueda rehacer la cuenta con el criterio que prefiera.

\subsection{Y la conclusión que el throttling deja clara}

El throttling confirma que los perfiles del manual \textbf{aprietan de verdad}:
hasta \num{2086}~s acumulados en una ventana de 300~s. Pero no es lo que
diferencia a los motores. Lo que los diferencia es si la consulta cabe en memoria:
donde cabe (\texttt{Q2}, por el \emph{top-K}), gana el motor con el plan más
eficiente; donde no cabe (\texttt{Q1} en PostgreSQL), el disco domina y la CPU
pasa a segundo plano.

% 07-recomendaciones.tex — Sección 7: recomendaciones técnicas, derivadas de los
% datos de este experimento, no de opiniones generales sobre los motores.

\section{Recomendaciones técnicas para despliegues documentales en la nube}
\label{sec:recomendaciones}

Las recomendaciones salen de lo medido aquí, con sus límites: un solo conjunto de
datos de \num{500000} documentos, dos consultas y un host virtualizado. No son
afirmaciones generales sobre los dos motores.

\paragraph{1. Elegir el motor por la forma de la consulta, no por el motor.}
Es el resultado más claro: MongoDB gana \texttt{Q1} en las nueve celdas y
PostgreSQL gana \texttt{Q2} en ocho de nueve. Si la carga dominante desanida
arreglos y agrupa, MongoDB rinde de 2 a 5 veces más; si ordena y se queda con un
Top-K, PostgreSQL rinde de 1.2 a 10.7 veces más. Un despliegue que haga las dos
cosas no se resuelve eligiendo motor, sino decidiendo qué consulta importa más.

\paragraph{2. Dimensionar \texttt{work\_mem} mirando el plan, no el tamaño de la tabla.}
\texttt{Q2} ordena \num{175049} filas y necesita \textbf{273~kB}, porque el
\texttt{LIMIT} la convierte en un \emph{top-K}. \texttt{Q1} sobre los mismos datos
necesita unos 0.8~GB \textbf{por consulta} y vuelca a disco con cualquier valor
razonable de \texttt{work\_mem}. Subir el parámetro de 4 a 64~MB no cambió nada
medible en \texttt{Q2} y no habría salvado a \texttt{Q1}: lo que decide es la forma
del plan.

\paragraph{3. En PostgreSQL, vigilar \texttt{temp\_bytes} como métrica de primera línea.}
El volcado no se nota en la CPU ni en el RSS: se nota en \textbf{\cifraTempTotalTB{}~TB}
escritos en disco a lo largo del experimento. Una consulta que vuelca degrada por
una vía que el throttling de CPU y la memoria residente no muestran, y
\texttt{pg\_stat\_database.temp\_bytes} la delata de inmediato.

\paragraph{4. No usar la memoria residente para dimensionar un contenedor de PostgreSQL.}
\texttt{max\_rss\_mb} no cuenta los \texttt{shared\_buffers}, que son memoria
compartida. En este experimento PostgreSQL aparece con 16~MiB residentes y un
conjunto de trabajo de 891~MiB en la misma celda. Un límite de memoria fijado con
la primera cifra mataría el contenedor de inmediato. La métrica útil es
\texttt{working\_set} (o \texttt{memory.current} menos la caché recuperable).

\paragraph{5. Con límites de CPU bajos, el calentamiento tiene que ser largo.}
Con 50 clientes y la caché limpia, la primera consulta de cada cliente tarda de
69 a 110~s en los perfiles B y C. Un calentamiento de 60~s deja ese costo dentro
de la ventana de medición y contamina los percentiles de cola: medido, hasta el
91~\% de las consultas de una ventana eran primeras consultas. Quien mida latencia
de cola bajo contención debe comprobar que su calentamiento cubre la primera
consulta, no suponerlo.

\paragraph{6. Un límite de memoria sin \emph{swap} convierte la saturación en caída, no en lentitud.}
Con \texttt{memory.swap.max = 0}, PostgreSQL con \texttt{Q1} y 50 clientes no se
degradó: fue \textbf{matado por el kernel}, y en tres de las cinco veces antes de
empezar a medir. Para un servicio en producción eso es la diferencia entre
responder despacio y no responder. Si el objetivo es degradar en vez de caer, hace
falta limitar la concurrencia en la capa de conexión antes de que el motor agote su
límite.

\paragraph{7. Medir en un portátil obliga a declarar la frecuencia del procesador.}
Durante la matriz el rendimiento del procesador osciló entre el 86~\% y el
60.8~\% de la frecuencia nominal, y la mayor dispersión observada coincide con una
de esas caídas. En un host dedicado esto no aparecería; en un entorno como este,
hay que medirlo y declararlo, no suponer que la CPU es constante.


\appendix
% A1-anexo-variantes.tex — Anexo: las dos variantes de estrés de memoria de Q2.
% Requisito R-16.1: cobertura parcial, decidida y justificada, FUERA del CSV de
% 36 celdas, que no se toca.

\section{Anexo: las dos variantes de estrés de memoria de \texttt{Q2}}
\label{sec:variantes-anexo}

\subsection{Por qué este anexo existe, y por qué es parcial}

El manual (§5.2) pide ejecutar \texttt{Q2} con dos configuraciones de memoria. El
conteo de escenarios que fija el propio manual son 36 celdas, y con 36 el esquema
del CSV \textbf{no tiene ninguna columna donde distinguir la variante}: añadir una
está prohibido explícitamente. Las dos exigencias no caben juntas.

La decisión tomada fue medir la segunda variante \textbf{en anexo}, fuera del CSV
de las 36 celdas, y solo en \textbf{\texttt{PROFILE\_C\_SEVERE}}: los dos motores y
las tres concurrencias, 6 celdas y 18 repeticiones. El razonamiento:

\begin{itemize}
  \item \texttt{PROFILE\_C} es el perfil con \textbf{menos memoria} (1~GiB). Si las
        dos variantes no se diferencian ahí, no se diferencian en ninguno.
  \item Por el plan ya se sabía que el ordenamiento de \texttt{Q2} ocupa
        \textbf{273~kB} (sección~\ref{sec:planes}), muy por debajo tanto de los
        4~MB de la variante restrictiva como de los 64~MB de la permisiva.
\end{itemize}

\textbf{La cobertura es parcial y se declara como tal}: no se midió
\texttt{Q2\_MEM\_HIGH} en los perfiles A y B. Los datos están en
\texttt{results/f8\_anexo/measurements.csv} (18 filas, validadas con el mismo
esquema de 14 columnas).

\subsection{Resultado: las dos variantes no se diferencian}

\begin{table}[htbp]
  \centering
  \caption{\texttt{Q2} en \texttt{PROFILE\_C\_SEVERE}: variante restrictiva
           (\texttt{work\_mem} 4~MB / \texttt{allowDiskUse} falso) frente a
           permisiva (64~MB / verdadero). Throughput medio de 3 repeticiones.}
  \label{tab:variantes}
  \small
  \begin{tabular}{lrrrrl}
    \toprule
    Motor & c & Restrictiva & Permisiva & Diferencia & Volcado \\
    \midrule
    MongoDB    & 1  & 0.0533 & 0.0444 & $-16.7$~\% & ninguno en ninguna \\
    MongoDB    & 10 & 0.1000 & 0.1000 & $0.0$~\%   & ninguno en ninguna \\
    MongoDB    & 50 & 0.1667 & 0.1667 & $0.0$~\%   & ninguno en ninguna \\
    PostgreSQL & 1  & 0.5722 & 0.6622 & $+15.7$~\% & ninguno en ninguna \\
    PostgreSQL & 10 & 0.6000 & 0.6222 & $+3.7$~\%  & ninguno en ninguna \\
    PostgreSQL & 50 & 0.5611 & 0.5744 & $+2.4$~\%  & ninguno en ninguna \\
    \bottomrule
  \end{tabular}
\end{table}

Las dos diferencias mayores (un $-16.7$~\% en MongoDB y un $+15.7$~\% en
PostgreSQL, las dos con un cliente) \textbf{apuntan en direcciones opuestas}, que
es la forma que tiene el ruido, no un efecto del parámetro. Las otras cuatro están
entre el 0 y el 3.7~\%, por debajo de la dispersión propia de esas celdas. Las 18
repeticiones salieron \texttt{ok} y \textbf{ninguna volcó a disco}, ni con
\texttt{allowDiskUse: true} en MongoDB ni con 4~MB de \texttt{work\_mem} en
PostgreSQL.

\subsection{Conclusión}

\textbf{La premisa del manual no se cumple.} Se esperaba que la variante
restrictiva forzara volcado a disco y que las dos configuraciones dieran resultados
distintos; medidas las dos en el perfil más exigente, no se diferencian y ninguna
vuelca. La causa está en el plan y no en el parámetro: el \texttt{LIMIT 1000}
acota el ordenamiento a 273~kB, así que el límite de memoria que el manual varía
nunca se alcanza. Se reporta como hallazgo, no como fallo del experimento.

% A2-limitaciones.tex — pendiente: se escribe con los datos ya generados.
%
% NINGÚN número a mano: usar \input de tablas/ y las macros de cifras/.



\end{document}
